% =============================================================================
%  SfN Neuroscience 2026 poster (Washington DC, Nov 14-18) -- paper Figs 1-5.
%  Board is 6 ft x 4 ft (72 x 48 in, SfN 2026 poster-session page); the poster is
%  60 x 40 in so it fits with a margin for pins and prints on a 42 in roll.
%  To change size, edit paperwidth/paperheight in the geometry line (type sizes are absolute).
%  Figures are the assembled vector PDFs, pulled straight from the paper:
%    Fig 1  -> Closedloop_nc/images/Figure1.pdf  (newest, with latency panel)
%    Fig 2-5 -> paper/figures_final/FigureN.pdf
%  Build:  latexmk -pdf poster.tex      (from this folder)
%  [FILL] marks what only the authors can supply.
% =============================================================================
\documentclass{article}
\usepackage[paperwidth=60in,paperheight=40in,margin=0.9in]{geometry}
\usepackage[T1]{fontenc}
\usepackage[scaled]{helvet}
\renewcommand{\familydefault}{\sfdefault}
\usepackage{amssymb}
\usepackage{sfmath}
\usepackage{anyfontsize}
\usepackage{graphicx}
\usepackage{xcolor}
\usepackage{enumitem}
\usepackage[most]{tcolorbox}
\pagestyle{empty}
\setlength{\parindent}{0pt}

\graphicspath{{../../paper/figures_final/}}
% Fig 1 is named by full path: figures_final/Figure1.pdf is a stale Fig-2 assembly.
\newcommand{\figonepath}{../../../Closedloop_nc/images/Figure1.pdf}

% ---- palette: navy header like the SfN template; CL blue / OL red from the figures
\definecolor{navy}{RGB}{31,58,95}
\definecolor{olred}{RGB}{220,60,45}
\definecolor{clblue}{RGB}{50,95,200}
\definecolor{ink}{RGB}{30,30,30}
\definecolor{fill}{RGB}{200,30,30}
\newcommand{\FILL}[1]{\textcolor{fill}{\textbf{[FILL: #1]}}}

% ---- type sizes (pt) -- body ~30 pt reads at 1.5 m
\newcommand{\bodysize}{\fontsize{28}{34}\selectfont}
\newcommand{\smallsize}{\fontsize{22}{27}\selectfont}
\newcommand{\headsize}{\fontsize{40}{48}\selectfont}

\setlist[itemize]{leftmargin=1.1em,itemsep=0.25em,topsep=0.3em,label=\textcolor{navy}{\(\blacktriangleright\)}}

% numbered result box: big number badge + claim as the heading
\newtcolorbox{resultbox}[2]{enhanced,colback=white,colframe=navy,boxrule=3pt,arc=0pt,
  left=14pt,right=14pt,top=10pt,bottom=12pt,toptitle=6pt,bottomtitle=6pt,
  title={\headsize\bfseries\begin{tabular}{@{}p{1.1em}@{}p{0.93\linewidth}@{}}\fontsize{56}{56}\selectfont #1 & \raggedright\hyphenpenalty=10000 #2\end{tabular}},
  coltitle=white,colbacktitle=navy,
  before upper={\bodysize\color{ink}}}
% plain section box (Question, Conclusions, Methods)
\newtcolorbox{plainbox}[1]{enhanced,colback=navy!6,colframe=navy,boxrule=3pt,arc=0pt,
  left=14pt,right=14pt,top=10pt,bottom=12pt,
  title={\headsize\bfseries #1},coltitle=white,colbacktitle=navy,
  before upper={\bodysize\color{ink}}}

\newcommand{\fig}[2][\linewidth]{\begin{center}\includegraphics[width=#1]{#2}\end{center}\vspace{-0.4em}}
\newcommand{\gap}{\vspace{0.4in}}

\begin{document}
% =============================================================================
%  HEADER
% =============================================================================
\begin{tcolorbox}[enhanced,colback=navy,colframe=navy,arc=0pt,boxrule=0pt,
   left=30pt,right=30pt,top=24pt,bottom=24pt,height=5in,valign=center]
\color{white}
\begin{minipage}[c]{0.12\linewidth}\centering
  \fontsize{26}{30}\selectfont \FILL{QR to preprint}
\end{minipage}\hfill
\begin{minipage}[c]{0.74\linewidth}\centering
  {\fontsize{92}{104}\selectfont\bfseries Closed-loop optogenetics steers mesoscale cortical activity\\[0.15em]
   and exposes the brain-state disturbances it must reject\par}
  \vspace{0.45em}
  {\fontsize{48}{56}\selectfont Aditya Deole, Anna Li, Eric Shea-Brown, Mehran Mesbahi, Nicholas A.\ Steinmetz\par}
  \vspace{0.2em}
  {\fontsize{38}{46}\selectfont University of Washington, Seattle, WA \quad \FILL{departments}\par}
\end{minipage}\hfill
\begin{minipage}[c]{0.12\linewidth}\centering
  {\fontsize{60}{64}\selectfont\bfseries \FILL{PSTR\#\#\#.\#\#}\par}
  \vspace{0.6em}
  {\fontsize{26}{30}\selectfont \FILL{UW logo}\par}
\end{minipage}
\end{tcolorbox}

\vspace{0.4in}
% =============================================================================
%  THREE COLUMNS
% =============================================================================
\noindent
\begin{minipage}[t]{0.32\linewidth}\vspace{0pt}
% ----------------------------------------------------------------------- col 1
\begin{plainbox}{Question}
Mesoscale cortical dynamics are linked to perception, cognition and movement, but
testing what they \emph{cause} requires shaping them precisely, not just perturbing them.
\begin{itemize}
  \item Can feedback from wide-field imaging hold one cortical region at a chosen level in real time?
  \item What limits it, and what does the error left under control tell us about brain state?
\end{itemize}
\end{plainbox}
\gap
\begin{resultbox}{1}{A real-time imaging + optogenetics control loop}
\fig[0.9\linewidth]{\figonepath}
\begin{itemize}
  \item Wide-field GCaMP imaging at 35\,Hz; the output is $\Delta F/F$ in a kernel beside a
        red-laser inhibition spot (a, b).
  \item Loop latency 24\,ms median, $\le$\,47\,ms -- inside one frame (c, d).
  \item Controller: PI feedback ($K_p, K_i$) plus a feedforward gain $K_r$ (e).
\end{itemize}
\end{resultbox}
\gap
\begin{resultbox}{2}{Trial-averaged responses are near-linear; single trials are not}
\fig{Figure2.pdf}
\begin{itemize}
  \item Suppression scales linearly with laser power (B; $R^2 = 0.73$--$0.97$, 4 sessions, 3 mice).
  \item A low-order LTI model predicts held-out trial averages (C, D; median $R^2 = 0.85$);
        the 0.11\,s dip timescale is shared across sessions (F).
  \item Models transfer only partly between sessions (E) $\Rightarrow$ use them for a rough
        feedforward gain and let feedback correct the rest.
  \item Single-trial variability follows state: $0.58\times$ in high motion, $2.35\times$ in high
        absolute $\delta$ power -- and stimulus-free activity does the same (G).
\end{itemize}
\end{resultbox}
\end{minipage}\hfill
% ----------------------------------------------------------------------- col 2
\begin{minipage}[t]{0.32\linewidth}\vspace{0pt}
\begin{resultbox}{3}{Feedback tracks the target better and cuts trial-to-trial variability}
\fig[0.96\linewidth]{Figure3.pdf}
\vspace{-0.2em}
{\smallsize Target $-5\%\,\Delta F/F$ for 3\,s; \textcolor{clblue}{\textbf{closed loop}} and
\textcolor{olred}{\textbf{open loop}} (feedforward only) randomly interleaved. 15 sessions, 4 mice.\par}
\begin{itemize}
  \item Tracking error (RMSE) $-23\%$ over the stimulus (ratio 0.77, $p = 5.8\times10^{-8}$),
        lower in \textbf{all 15 sessions}; $-32\%$ at steady state (H).
  \item Across-trial variance $-42\%$ (0.58, $p = 6.3\times10^{-6}$); the ratio is 1 before and
        after stimulation $\Rightarrow$ \textbf{disturbance rejection}, not trial selection (D, E, I).
  \item Still 13\% better than the best \emph{fixed} input chosen in hindsight ($p = 6.5\times10^{-4}$),
        so the gain is feedback, not calibration.
\end{itemize}
\end{resultbox}
\gap
\begin{resultbox}{4}{Performance depends on brain state; feedback rejects a disturbance predicted from the other hemisphere}
\fig[0.9\linewidth]{Figure4.pdf}
\begin{itemize}
  \item Transient error is set by the initial deviation (unique $R^2 = 0.53$); steady-state error
        by slow, large absolute 1--4\,Hz power (B, C).
  \item The more the mouse moves, the larger the share of error feedback removes
        ($-13\%$ per SD, $p = 0.0023$); high absolute $\delta$ power raises closed-loop error
        ($p = 7.1\times10^{-4}$) (D).
  \item A contralateral$\to$ipsilateral model of spontaneous activity (held-out $R^2 = 0.76$)
        estimates each trial's disturbance (E, F). Closed loop leaves 0.58 of its energy
        (rejection ratio, $p = 1.3\times10^{-4}$, 11 sessions) (G).
\end{itemize}
\end{resultbox}
\end{minipage}\hfill
% ----------------------------------------------------------------------- col 3
\begin{minipage}[t]{0.32\linewidth}\vspace{0pt}
\begin{resultbox}{5}{Tracking a moving target: feedback cuts error, preview cancels lag}
\fig{Figure5.pdf}
{\smallsize 1\,Hz sinusoidal target; preview shifts the reference ahead by the cortical delay
(5 samples, 143\,ms). OL, OL$+$preview, CL, CL$+$preview interleaved (Mouse 3, 3 sessions).\par}
\begin{itemize}
  \item Feedback halves across-trial variance (11.8 $\to$ 5.8\,$(\%\Delta F/F)^2$) and lowers
        median RMSE 3.59 $\to$ 2.26 (E--G).
  \item Preview cancels the phase lag (OL $71^\circ \to 16^\circ$, CL $46^\circ \to 2^\circ$),
        as the lookahead predicts (H).
  \item Same direction in all 3 sessions (I--K; $n = 3$, no tests).
\end{itemize}
\end{resultbox}
\gap
\begin{plainbox}{Conclusions}
\begin{itemize}
  \item Approximately linear, time-invariant trial-averaged dynamics are enough to design a
        controller that \textbf{holds regional cortical activity at a target}.
  \item Feedback beats open loop in every session and mouse, and \textbf{rejects ongoing disturbances}
        -- part of them predictable from the opposite hemisphere.
  \item The error left under control is \textbf{a readout of brain state}: it grows with slow
        $\delta$-band fluctuations, while feedback removes more of the error that comes with movement.
  \item Feedback and preview fix different failures; together they track a moving target.
\end{itemize}
\end{plainbox}
\gap
\begin{plainbox}{Methods}
{\smallsize
CaMKII-tTA;tetO-GCaMP8s mice; systemic AAV PHP.eB DLX2.0-ChrimsonR (inhibitory interneurons),
so 638\,nm light suppresses local excitatory activity. Head-fixed, awake, no task.
Wide-field imaging (470/405\,nm alternating, 35\,Hz); galvo-steered laser.
Closed-loop $\Delta F/F$ uses a trailing 40\,s mean (the loop's own online baseline).
Statistics: per-session log CL/OL ratios, two-sided one-sample $t$-test across sessions,
checked per mouse. Rig and preparation: Matveev et al.\ (2024).\par}
\end{plainbox}
\gap
{\smallsize\color{ink}
\textbf{Funding:} \FILL{grants -- SfN requires a funding acknowledgement}\\[0.3em]
\textbf{Contact:} \FILL{email}\par}
\end{minipage}

\end{document}
